% =========================================================================
% sec01_introduction.tex — Section I of the sealed master
% sealed master paper_seni142_ieee_transactions_en.md; editorial injection pass 2026-10-17.
% =========================================================================
\section{Introduction}
\label{sec:intro}

Insular interconnected systems occupy an extreme point of the power engineering design space. An island cannot draw inertia, reactive support, or emergency transfer from neighbors: every megawatt of reserve, every volt-ampere of reactive compensation, and every megahertz of damping capability must be produced, dispatched, and defended inside the same synchronous footprint that consumes it. The Dominican Republic is a canonical instance: a single synchronous island of roughly eleven million people, with no submarine or overland interconnection, more than forty percent of national load concentrated in the metropolitan Ozama region, an inertia anchor dominated by the Punta Catalina coal plant, and a renewable build-out whose instantaneous solar share has reached 36.6\% of demand in operational telemetry \cite{ocseni}. The national consequences of losing the generation--load equilibrium are not hypothetical: the nationwide service interruptions recorded in November 2025 and February 2026, documented in operator event logs \cite{ocseni} and in the national press archive \cite{listin2025}, mobilized emergency response at the scale of a full territory and renewed the policy debate on frequency stability, black-start capability, and the observability of the transmission estate \cite{blackout2004}. The physics of this fragility is by now classical \cite{kundur1994} --- declining system inertia degrades the rate of change of frequency and narrows the feasible control envelope of converter-dominated grids \cite{tielens2016}, \cite{milano2018}, \cite{dorfler2013}, \cite{hatziargyriou2021} --- but the empirical instruments available to study it remain surprisingly primitive.

The benchmarking infrastructure of the field is, in its majority, synthetic. The canonical test cases that still govern doctoral training and comparative studies --- the WSCC 9-bus network, the IEEE 14-, 30-, 57-, and 118-bus systems --- descend from drafting-table abstractions of the early 1960s: reactances are round academic numbers, loads are fixed triangles of power, geography is absent, climate is absent, and nothing in the dataset has ever tripped, been maintained, been paid for, or been hit by a hurricane. At the opposite pole, utility-grade network cases are held in proprietary formats of commercial simulation platforms, are encumbered by confidentiality, and are rarely released in forms that permit independent re-execution. The open-systems movement has answered with valuable infrastructure --- the MATPOWER case collection \cite{matpower2011}, and the Power Grid Library of AC-OPF benchmark instances \cite{pglib2019} among them --- but its instances are, almost without exception, \emph{synthetic or anonymized}: they are grids that exist to be hard, not grids that exist. Between the toy and the secret there is almost nothing: an open, complete, measured model of a real national transmission system, with provenance sufficient for independent verification, is a public good that the research community currently lacks for the insular Caribbean --- and, with few exceptions, for the insular world at large.

This article closes that gap with \textbf{SENI-142}, an open-access benchmark constructed around three commitments. First, \textbf{fidelity to a measured reality}: the network model is not synthetic --- it is derived from the national transmission estate and validated against exactly 520 complete consecutive 168-hour weeks --- 87,360 hourly states spanning 3,640 Monday-to-Sunday days, drawn without a single interpolated hour from the 2016--2025 operational archive of the system operator --- dispatched unit by unit and hour by hour, so that every simulated state answers to an archived operational one. Second, \textbf{structural honesty about scale}: the full national system is solved --- all bulk buses, all trunk lines, all transformers --- not a reduced district of it, and the solution method is formulated so that its exactness and complexity are theorems rather than aspirations: a hierarchical border-coordinate formulation (HBCF) in the diakoptic tradition \cite{kron1963}, \cite{happ1971}, condensed by the Schur complement \cite{zhang2005} and factorized on a partition that follows the ten territorial planning regions mandated by national spatial planning law (Law 368-22) \cite{ley368}. Third, \textbf{deterministic reproducibility}: every artifact passes a four-phase digital-twin lifecycle with a fail-closed gate cascade, detached SHA-512 sealing, and atomic promotion, so that any figure reported in this article can be regenerated from a bit-identical, cryptographically attested software state.

The contributions are five, and each is falsifiable:
\begin{itemize}
\item \textbf{C1 --- Benchmark.} An open, georeferenced, standards-consistent model of the complete bulk transmission network of the Dominican Republic: 142 nodal buses over approximately 112 physical substations, 135 transmission lines, 58 power transformers, strata of 345, 230, and 138~kV, with the 69~kV subtransmission layer reduced by injection-preserving equivalents; released with topology, parameters, provenance metadata, and validation harness through IEEE DataPort and GitHub.
\item \textbf{C2 --- Formulation.} The HBCF diakoptic partition of a full insular grid into ten sovereign territorial microgrids --- the legal planning regions of Law 368-22 --- aggregated into four electrical macro-subdomains with a minimal border set, and the proof that its Schur condensation is exact, with complexity $O(\sum_k N_k^3 + N_B^3)$ against the $O(N^3)$ of monolithic dense factorization.
\item \textbf{C3 --- Deterministic digital twin.} A four-phase lifecycle (hermetic POSIX-namespace enclave with a fourteen-gate fail-closed cascade; research-testbed verification of static post-event $N-1$ feasibility against a 215-contingency catalog; SHA-512-detached Golden Master; atomic kernel promotion via \texttt{renameat2} with zero regression) that converts simulation telemetry into evidence with a chain of custody.
\item \textbf{C4 --- Multi-vector flexibility.} The coupling of power-to-gas hydrogen (electrolysis, 700~bar storage, fuel-cell re-injection), 764~MW of reversible pumped hydro storage, and ADM1-based biomethanation to a multi-temporal AC-OPF, stated with its full Karush--Kuhn--Tucker (KKT) system and validated by a zonal CAPEX/OPEX balancing proof.
\item \textbf{C5 --- Measured silicon performance.} Resolution of each hourly national AC state in 2.65~ms with 473~MiB of resident memory and fasorial border residual $\delta = 6.35 \times 10^{-12}\ \mathrm{pu}$ on consumer-grade hardware --- two to three orders of magnitude inside the 1--5~s state-estimation cycle of an energy management system (EMS), adequate for static contingency screening --- under a 24-hour block-bootstrap validation protocol.
\end{itemize}

The remainder of the article is organized as follows. Section~\ref{sec:methodology} states the methodology: the invariant formulation (graph model, 69~kV reduction, HBCF partition, Schur condensation, data protocol, deterministic lifecycle, reactive-limit regularization, multi-temporal AC-OPF) and the modular clause that instantiates the Dominican case under Law 368-22. Section~\ref{sec:topology} documents the benchmark topology, from the 112-substation physical estate to the 142-bus nodal model and the ten territorial microgrids. Section~\ref{sec:multivector} develops the multi-vector flexibility constellation and its zonal economic validation. Section~\ref{sec:benchmarking} reports silicon performance, accuracy, and contingency resilience. Section~\ref{sec:discussion} discusses policy implications and limitations, and Section~\ref{sec:conclusion} concludes with the open-release plan and four research routes.
